\section{Forelesning 4: \emph{Strategi og digital transformasjon}}
\label{sec:lec04}

\subsection*{Oversikt}
Forelesning 4 (Ulrik Røhl) gjør et avgjørende begrepsmessig skille mellom
\term{Digital Transformation} og \term{IT-Enabled Organizational Transformation},
basert på Wessel et al.\ (2021). Distinksjonen er strategisk viktig fordi den
bestemmer \emph{hvilken type} digital investering som faktisk passer
virksomheten -- og dermed hvilken anbefaling som er forsvarlig.

\subsection{Hovedteori}

\subsubsection*{Wessel et al.\ (2021): De to typene transformasjon}

\begin{definition}{Digital Transformation (DT)}
Digital teknologi er \emph{sentral for å redefinere} virksomhetens value
propositions, noe som fører til en \emph{ny organisasjonsidentitet}.
Virksomheten blir noe annet enn den var.
\end{definition}

\begin{definition}{IT-Enabled Organizational Transformation (ITEOT)}
Digital teknologi \emph{støtter og forsterker en eksisterende} value
proposition og identitet. Virksomheten gjør det samme den alltid har gjort,
men bedre, raskere eller billigere.
\end{definition}

\keypoint{Forskjellen ligger i hva som skjer med \emph{organisasjonsidentiteten}:
endres den (DT) eller forsterkes den (ITEOT)?}

\subsubsection*{Empirisk grunnlag (Wessel et al.\ 2021)}
\begin{itemize}
  \item \textbf{Case Alpha} -- Et sykehus i Frankrike som innfører elektronisk
    pasientjournal. Identiteten forblir «sykehus». Klassisk
    \emph{IT-enabled transformation}.
  \item \textbf{Case Beta} -- En produksjonsbedrift i Finland som beveger seg
    fra maskinvare til programvare-kontroll. Sitatet er sterkt: «Vi var
    maskinfabrikanter \ldots\ men nå trenger ikke kontrollprogramvaren
    maskinvare lenger.» Identiteten endres fundamentalt -- ekte
    \emph{digital transformation}.
\end{itemize}

\subsubsection*{Hvorfor distinksjonen er strategisk avgjørende}
Feiltolkning fører til feilinvestering:
\begin{itemize}
  \item \emph{Tror} man at virksomheten gjennomgår digital transformasjon, kan
    man argumentere for å bygge en full plattformidentitet -- med
    tilsvarende risiko.
  \item \emph{Ser} man korrekt at det er IT-enabled transformasjon, fokuseres
    investeringen på \emph{hvor digitalt støtter} det eksisterende
    verdiforslaget.
\end{itemize}

\subsubsection*{Kobling til posisjonering (\autoref{sec:lec01})}
DT redefinerer hva virksomheten \emph{er} -- en ny strategisk posisjon. ITEOT
forsterker den eksisterende posisjonen. Porters spørsmål («er aktiviteten en
del av et unikt aktivitetssystem?») gjelder begge, men svaret ser ulikt ut.

\subsection{Sentrale fagbegreper}
\begin{itemize}
  \item \term{Digital Transformation} -- redefinering av value proposition og
    identitet.
  \item \term{IT-Enabled Organizational Transformation} -- forsterkning av
    eksisterende value proposition.
  \item \term{Organizational identity} -- selvforståelsen av «hva vi er som
    virksomhet».
  \item \term{Value proposition} -- løftet til kunden.
\end{itemize}

\subsection{Modeller og rammeverk}

\figureplaceholder{lec04-dt-vs-iteot}{Wessel et al.\ (2021): To typer
transformasjon -- DT vs.\ ITEOT, med identitetsendring som skillemarkør.}

\subsubsection*{Diagnostisk spørsmål}
For å plassere en konkret virksomhet, still:
\begin{enumerate}
  \item Endres value proposition fundamentalt eller forsterkes den eksisterende?
  \item Endres organisasjonsidentiteten -- blir virksomheten «noe annet»?
  \item Er den digitale teknologien \emph{sentral} (uten den finnes ikke
    tilbudet) eller \emph{støttende} (tilbudet finnes uten, men forbedres)?
\end{enumerate}

\subsection{Eksempler og caser}
\begin{itemize}
  \item \textbf{Coop} gjennomgår \emph{IT-enabled transformation}. Appen
    støtter og forsterker den butikkbaserte forretningsmodellen og
    medlemsidentiteten. Coop redefinerer seg ikke som digital aktør --
    bevis: lukkingen av Coop.dk-nettbutikken i januar 2025.
  \item \textbf{Netflix} -- klassisk DT: gikk fra DVD-utleie til
    streaming-/produksjonsselskap. Identiteten endret seg.
  \item \textbf{Finlandsk maskinfabrikant (Beta)} -- DT: fra maskinvare til
    programvarekontroll.
  \item \textbf{Banker som lanserer mobilbank} -- ITEOT så lenge identiteten
    forblir «bank»; ville blitt DT hvis banken redefinerte seg som
    plattformaktør i bredere finans-/dataøkosystem.
\end{itemize}

\subsection{Eksamensrelevans}
\begin{itemize}
  \item Kunne definere begge typer presist (Wessel et al.\ 2021).
  \item Bruke distinksjonen til å diagnostisere konkrete virksomheter.
  \item Forsvare \emph{hvorfor} distinksjonen er strategisk viktig (forhindrer
    feilinvestering).
  \item Kombinere med Weill \& Woerners DBM (\autoref{sec:lec03}): Coop som
    \emph{omnichannel} + ITEOT $\rightarrow$ digital investering bør styrke
    butikkøkonomien, ikke bygge en separat digital identitet.
\end{itemize}
